\frfilename{mt216.tex}
\versiondate{25.9.04}
\copyrightdate{1994}

\def\chaptername{Taksonomi ruang ukur}
\def\sectionname{Contoh-contoh}

\newsection{216}

Sudah menjadi kebiasaan umum -- dan, menurut pandangan saya, kebiasaan
yang baik -- dalam buku-buku matematika murni, untuk memberikan contoh
pembeda;  maksud saya, setiap kali kita diberi suatu daftar konsep baru,
kita berharap akan diberi contoh-contoh yang menunjukkan bahwa kita
memiliki gambaran yang memadai tentang hubungan di antara konsep-konsep
tersebut, dan khususnya bahwa tidak ada implikasi mengejutkan yang
disembunyikan dari kita.   Mengenai konsep-konsep yang didaftar dalam
211A-211K, %211A 211B 211C 211D 211E 211F 211G 211H 211I 211J 211K
kita memiliki sepuluh sifat berbeda yang dimiliki oleh sebagian, tetapi
tidak semua, ruang ukur, sehingga secara teoretis terdapat total $2^{10}$
jenis ruang ukur yang berbeda, yang diklasifikasikan menurut sifat mana
di antara sepuluh sifat tersebut yang dimilikinya.   Daftar hubungan dasar
dalam 211L mengurangi 1024 kemungkinan ini menjadi 72.   Dengan mengamati
bahwa suatu ruang dapat sekaligus tanpa atom dan atomik murni hanya jika
ukuran seluruh ruang adalah $0$, kita memperoleh 56 kemungkinan, yaitu dua
kasus sepele dengan $\mu X=0$ (karena ukuran seperti itu dapat lengkap atau
tidak lengkap) bersama dengan $9\times 2\times 3$ kasus, yang bersesuaian
dengan sembilan kelas

\quad ruang peluang,

\quad ruang yang berhingga total, tetapi bukan ruang peluang,

\quad ruang yang $\sigma$-hingga, tetapi tidak berhingga total,

\quad ruang yang dapat dilokalkan secara ketat, tetapi tidak $\sigma$-hingga,

\quad ruang yang dapat dilokalkan dan ditentukan secara lokal, tetapi tidak
dapat dilokalkan secara ketat,

\quad ruang yang dapat dilokalkan, tetapi tidak ditentukan secara lokal,

\quad ruang yang ditentukan secara lokal, tetapi tidak dapat dilokalkan,

\quad ruang yang semihingga, tetapi tidak ditentukan secara lokal maupun
dapat dilokalkan,

\quad ruang yang tidak semihingga;

\noindent dua kelas

\quad ruang yang lengkap,

\quad ruang yang tidak lengkap;

\noindent dan tiga kelas

\quad ruang yang tanpa atom, dengan ukuran bukan $0$,

\quad ruang yang atomik murni, dengan ukuran bukan $0$,

\quad ruang yang tidak tanpa atom maupun atomik murni.

Saya tidak bermaksud memberikan seperangkat lengkap lima puluh enam
contoh, terutama karena gagasan yang diperlukan sebenarnya jauh
lebih sedikit daripada lima puluh enam.   Namun, saya sungguh berpendapat
bahwa untuk memahami ruang ukur abstrak dengan baik, kita perlu melihat
realisasi beberapa kombinasi sifat yang penting.   Karena itu, dalam
beberapa paragraf saya menjelaskan tiga contoh khusus sebagai tambahan
bagi contoh-contoh dalam 211M-211R. %211M 211N 211O 211P 211Q 211R

\leader{216A}{Ukuran Lebesgue}\cmmnt{ Sebelum beralih ke gagasan-gagasan baru, izinkan saya menyebut ukuran Lebesgue sekali lagi.
Sebagaimana dicatat dalam 211M, 211P dan 211Qa,

}(a) ukuran Lebesgue $\mu$ pada $\Bbb R$ lengkap, tanpa atom, dan
$\sigma$-hingga, sehingga dapat dilokalkan secara ketat, dapat
dilokalkan, dan ditentukan secara lokal.

(b) Ukuran subruang $\mu_{[0,1]}$ pada $[0,1]$ merupakan ukuran peluang
yang lengkap dan tanpa atom.

(c) Restriksi $\mu\restr\Cal B$ dari $\mu$ pada aljabar-sigma Borel
$\Cal B$ dari $\Bbb R$ bersifat tanpa atom, $\sigma$-hingga, dan tidak
lengkap.

\cmmnt{
\leader{216B}{}Sekarang saya mulai menjelaskan tiga `contoh tandingan';
yakni ruang-ruang yang dibangun secara khusus untuk menunjukkan bahwa
tidak ada implikasi tak terduga di antara sepuluh sifat yang sedang
dipertimbangkan di sini.   Bahkan menurut standar bab ini, contoh-contoh
tersebut harus dipandang sebagai bagian yang boleh dilewati oleh pelajar
yang ingin melanjutkan pekerjaan utama memahami teorema-teorema besar
dalam bidang ini.   Baik keberadaan contoh-contoh ini maupun teknik yang
diperlukan untuk membangunnya tidak mutlak diperlukan bagi hal lain yang
akan kita kaji sebelum Jilid 5.   Tetapi jika Anda hendak mempelajari teori
ukuran abstrak secara serius, cepat atau lambat Anda perlu membentuk suatu
gambaran mental tentang hakikat ruang-ruang yang memiliki berbagai sifat
di sini, dan syarat minimal bagi gambaran demikian ialah bahwa gambaran
itu mencakup perbedaan-perbedaan yang diperlihatkan oleh contoh-contoh ini.
}%end of comment

\leader{*216C}{Ruang lengkap yang dapat dilokalkan tetapi tidak ditentukan
secara lokal}\cmmnt{ Contoh pertama hampir tidak memerlukan gagasan di luar
yang sudah kita miliki, tetapi memang menuntut lebih banyak manipulasi
daripada yang layak diberikan sebagai latihan, sehingga mungkin berguna
sebagai peragaan teknik.

\medskip

}{\bf (a)} Misalkan $I$ sembarang himpunan tak terhitung, dan tetapkan
$X=\{0,1\}\times I$. Untuk $E\subseteq X$, $y\in\{0,1\}$ tetapkan
$E[\{y\}]=\{i:(y,i)\in E\}\subseteq I$.   Tetapkan

\Centerline{$\Sigma=\{E:E\subseteq X,\,E[\{0\}]\symmdiff E[\{1\}]$ terhitung$\}$.}

\noindent Maka $\Sigma$ merupakan aljabar-sigma dari himpunan-himpunan
bagian $X$.
\prooflet{\Prf\ (i)
$\emptyset[\{0\}]\symmdiff\emptyset[\{1\}]=\emptyset$ terhitung, jadi
$\emptyset\in\Sigma$.   (ii) Jika $E\in\Sigma$ maka

\Centerline{$(X\setminus E)[\{0\}]
\symmdiff(X\setminus E)[\{1\}]=E[\{0\}]\symmdiff E[\{1\}]$}

\noindent terhitung.   (iii) Jika $\sequencen{E_n}$ merupakan suatu
barisan dalam $\Sigma$ dan $E=\bigcup_{n\in\Bbb N}E_n$, maka

\Centerline{$E[\{0\}]\symmdiff E[\{1\}]
\subseteq\bigcup_{n\in\Bbb N}(E_n[\{0\}]\symmdiff E_n[\{1\}])$}

\noindent terhitung.\ \Qed}

Untuk $E\in\Sigma$, tetapkan $\mu E=\#(E[\{0\}])$ jika ini hingga,
$\infty$ jika tidak;  maka $(X,\Sigma,\mu)$ merupakan ruang ukur.

\medskip

{\bf (b)} $(X,\Sigma,\mu)$ lengkap.   \prooflet{\Prf\ Jika $A\subseteq
E\in\Sigma$ dan $\mu E=0$, maka $(0,i)\notin E$ untuk setiap $i$.   Jadi

\Centerline{$A[\{0\}]\symmdiff A[\{1\}]=A[\{1\}]\subseteq E[\{1\}]
=E[\{1\}]\symmdiff E[\{0\}]$}

\noindent harus terhitung, dan $A\in\Sigma$.\ \Qed}

\medskip

{\bf (c)} $(X,\Sigma,\mu)$ semihingga.   \prooflet{\Prf\ Jika $E\in\Sigma$
dan $\mu E>0$, terdapat $i\in I$ sedemikian sehingga $(0,i)\in E$;  kini
$F=\{(0,i)\}\subseteq E$ dan $\mu F=1$.\ \Qed}

\medskip

{\bf (d)} $(X,\Sigma,\mu)$ dapat dilokalkan.   \prooflet{\Prf\ Misalkan
$\Cal E$ sembarang himpunan bagian dari $\Sigma$.   Tetapkan

\Centerline{$J=\bigcup_{E\in\Cal E}E[\{0\}]$,
\quad $G=\{0,1\}\times J$.}

\noindent Maka $G\in\Sigma$.   Jika $H\in\Sigma$, maka

$$\eqalign{\mu(E\setminus H)&=0 \text{ untuk setiap }E\in\Cal E\cr
&\iff E[\{0\}]\subseteq H[\{0\}]\text{ untuk setiap }E\in\Cal E\cr
&\iff (0,i)\in H\text{ untuk setiap }i\in J\cr
&\iff \mu(G\setminus H)=0.\cr}$$

\noindent Jadi $G$ merupakan supremum esensial bagi $\Cal E$ dalam
$\Sigma$;  karena $\Cal E$ sembarang, $\mu$ dapat dilokalkan.\ \Qed}

\medskip

{\bf (e)} $(X,\Sigma,\mu)$ tidak ditentukan secara lokal.
\prooflet{\Prf\
Tinjau $H=\{0\}\times I$.   Maka $H\notin\Sigma$ karena
$H[\{0\}]\symmdiff H[\{1\}]=I$ tak terhitung.   Tetapi misalkan
$E\in\Sigma$ sembarang himpunan sedemikian sehingga $\mu E<\infty$.   Maka

\Centerline{$(E\cap H)[\{0\}]\symmdiff (E\cap H)[\{1\}]
=(E\cap H)[\{0\}]\subseteq E[\{0\}]$}

\noindent hingga, sehingga $E\cap H\in\Sigma$.   Karena $E$ sembarang,
$H$ menjadi saksi bahwa $\mu$ tidak ditentukan secara lokal.\ \Qed}
\medskip

{\bf (f)} $(X,\Sigma,\mu)$ atomik murni.   \prooflet{\Prf\ Misalkan
$E\in\Sigma$ sembarang himpunan dengan ukuran bukan nol.   Ambil $i\in I$
sedemikian sehingga $(0,i)\in E$.   Maka $(0,i)\in E$ dan
$F=\{(0,i)\}$ merupakan himpunan berukuran $1$ yang termuat dalam $E$;
karena $F$ merupakan himpunan tunggal, ia harus menjadi suatu atom bagi
$\mu$;  karena $E$ sembarang, $\mu$ atomik murni.\ \Qed}

\cmmnt{\medskip

{\bf (g)} Jadi konstruksi di sini menghasilkan ruang lengkap, dapat
dilokalkan, atomik murni, dan tidak ditentukan secara lokal.
}%end of comment

\leader{*216D}{Ruang lengkap yang ditentukan secara lokal tetapi tidak
dapat dilokalkan}\cmmnt{ Konstruksi berikut memerlukan sedikit teori
himpunan.}
Kita memerlukan dua himpunan $I$, $J$ sedemikian sehingga $I$ tak terhitung\cmmnt{ (lebih tepatnya, $I$ tidak dapat dinyatakan sebagai
gabungan terhitung dari himpunan-himpunan terhitung)},
$I\subseteq J$ dan $J$ tidak dapat dinyatakan sebagai
$\bigcup_{i\in I}K_i$ dengan setiap $K_i$ terhitung.   \cmmnt{Cara paling
alami untuk melakukan ini, dengan mengasumsikan aksioma pilihan, ialah
mengambil $I=\omega_1$, ordinal tak terhitung pertama, dan $J$ sebagai
$\omega_2$, ordinal pertama yang darinya tidak ada injeksi ke dalam $\omega_1$
(lihat 2A1Fc);  tetapi jika Anda lebih menyukai formulasi lain
(misalnya, $I=\{\{x\}:x\in\Bbb R\}$ dan $J=\Cal P\Bbb R$), saya akan
menuliskan argumen berikut dalam istilah $I$ dan $J$, dan Anda dapat
memilih pasangan Anda sendiri.}%end of comment

\medskip

{\bf (a)} Misalkan $\Tau$ adalah aljabar-sigma terhitung-koterhitung
dari $J$ dan $\nu$ ukuran terhitung-koterhitung pada
$J$\cmmnt{ (211R)}.   Tetapkan $X=J\times J$ dan untuk $E\subseteq X$
tetapkan

\Centerline{$E[\{\xi\}]=\{\eta:(\xi,\eta)\in E\}$,\quad
$E^{-1}[\{\xi\}]=\{\eta:(\eta,\xi)\in E\}$}

\noindent untuk setiap $\xi\in J$.   Tetapkan

\Centerline{$\Sigma=\{E:E[\{\xi\}]$ dan $E^{-1}[\{\xi\}]$ termasuk dalam
$\Tau$ untuk setiap $\xi\in J\}$,}

\Centerline{$\mu E=\sum_{\xi\in J}\nu E[\{\xi\}]
+\sum_{\xi\in J}\nu E^{-1}[\{\xi\}]$}

\noindent untuk setiap $E\in\Sigma$.   \cmmnt{Mudah diperiksa bahwa}
$\Sigma$ merupakan aljabar-sigma dan\cmmnt{ bahwa} $\mu$ merupakan
suatu ukuran.

\medskip

{\bf (b)} $(X,\Sigma,\mu)$ lengkap.   \prooflet{\Prf\ Jika $A\subseteq
E\in\Sigma$ dan $\mu E=0$, maka semua himpunan $E[\{\xi\}]$ dan
$E^{-1}[\{\xi\}]$ terhitung, sehingga hal yang sama berlaku bagi semua
himpunan $A[\{\xi\}]$ dan $A^{-1}[\{\xi\}]$, dan $A\in\Sigma$.\ \Qed}

\medskip

{\bf (d)} $(X,\Sigma,\mu)$ semihingga.   \prooflet{\Prf\ Untuk setiap
$\zeta\in J$, tetapkan

\Centerline{$G_{\zeta}=\{\zeta\}\times J$,\quad
$\tilde G_{\zeta}=J\times\{\zeta\}$.}

\noindent Maka semua penampang $G_{\zeta}[\{\xi\}]$,
$G_{\zeta}^{-1}[\{\xi\}]$, $\tilde G_{\zeta}[\{\xi\}]$, dan
$\tilde  G_{\zeta}^{-1}[\{\xi\}]$ adalah salah satu dari $J$,
$\emptyset$, atau $\{\zeta\}$, sehingga termasuk dalam $\Tau$, dan semua
$G_{\zeta}$, $\tilde G_{\zeta}$ termasuk dalam $\Sigma$, dengan ukuran
$\mu$ sebesar $1$.

Andaikan $E\in\Sigma$ merupakan himpunan dengan ukuran yang benar-benar positif.
Maka harus ada suatu $\xi\in J$ sedemikian sehingga

\Centerline{$0<\nu E[\{\xi\}]+\nu E^{-1}[\{\xi\}]
=\mu(E\cap G_{\xi})+\mu(E\cap \tilde G_{\xi})<\infty$,}

\noindent dan salah satu himpunan $E\cap G_{\xi}$,
$E\cap\tilde G_{\xi}$ merupakan himpunan berukuran hingga bukan nol yang
termuat dalam $E$.\ \Qed}

\medskip

{\bf (e)} $(X,\Sigma,\mu)$ ditentukan secara lokal.
\prooflet{\Prf\ Andaikan $H\subseteq X$ sedemikian sehingga
$H\cap E\in\Sigma$ setiap kali $E\in\Sigma$ dan $\mu E<\infty$.   Maka,
khususnya, $H\cap
G_{\zeta}$ dan $H\cap \tilde G_{\zeta}$ termasuk dalam
$\Sigma$, sehingga

\Centerline{$H[\{\zeta\}]
=(H\cap \tilde G_{\zeta})[\{\zeta\}]\in\Tau$,}

\Centerline{$H^{-1}[\{\zeta\}]
=(H\cap G_{\zeta})^{-1}[\{\zeta\}]\in\Tau$,}

\noindent untuk setiap $\zeta\in J$.   Ini menunjukkan bahwa
$H\in\Sigma$.   Karena $H$ sembarang, $\mu$ ditentukan secara lokal.\ \Qed}

\medskip

{\bf (f)} $(X,\Sigma,\mu)$ tidak dapat dilokalkan.
\prooflet{\Prf\ Tetapkan $\Cal E=\{G_{\zeta}:\zeta\in J\}$.
\Quer\ Andaikan, jika mungkin, bahwa $G\in\Sigma$ merupakan supremum
esensial bagi $\Cal E$.   Maka

\Centerline{$\nu(J\setminus G[\{\xi\}])
=\mu(G_{\xi}\setminus G)=0$}

\noindent dan $J\setminus G[\{\xi\}]$ terhitung, untuk setiap $\xi\in
J$.   Akibatnya $J\ne\bigcup_{\xi\in I}(J\setminus G[\{\xi\}])$, dan
terdapat suatu $\eta$ yang termasuk dalam
$J\setminus\bigcup_{\xi\in I}(J\setminus
G[\{\xi\}])=\bigcap_{\xi\in I}G[\{\xi\}]$. Ini berarti tepat bahwa
$(\xi,\eta)\in G$ untuk setiap $\xi\in I$, yaitu
$I\subseteq G^{-1}[\{\eta\}]$.   Dengan demikian
$G^{-1}[\{\eta\}]$ tak terhitung, sehingga
$\nu G^{-1}[\{\eta\}]=\mu(G\cap\tilde G_{\eta})=1$.   Tetapi perhatikan
bahwa $\mu(G_{\xi}\cap \tilde G_{\eta})=\mu\{(\xi,\eta)\}=0$ untuk
setiap $\xi\in
J$.   Ini berarti bahwa, dengan menetapkan
$H=X\setminus \tilde G_{\eta}$, $E\setminus H$ terabaikan untuk setiap
$E\in\Cal E$;  sehingga kita harus mempunyai
$0=\mu(G\setminus H)=\mu(G\cap\tilde G_{\eta})=1$, yang mustahil.\ \Bang

Jadi $\Cal E$ tidak memiliki supremum esensial dalam $\Sigma$, dan $\mu$
tidak dapat dilokalkan.\ \Qed}
\medskip

{\bf (g)} $(X,\Sigma,\mu)$ atomik murni.   \prooflet{\Prf\ Jika
$E\in\Sigma$ berukuran bukan nol, harus ada suatu $\xi\in J$ sedemikian
sehingga salah satu dari $E[\{\xi\}]$, $E^{-1}[\{\xi\}]$ tak terhitung;  yakni, sedemikian sehingga salah satu dari $E\cap G_{\xi}$,
$E\cap \tilde G_{\xi}$ tidak terabaikan.   Tetapi jika sekarang
$H\in\Sigma$ dan $H\subseteq E\cap G_{\xi}$, maka
$H[\{\xi\}]$ terhitung dan $\mu H=0$, atau
$J\setminus H[\{\xi\}]$ terhitung dan
$\mu(G_{\xi}\setminus H)=0$;  demikian pula, jika
$H\subseteq E\cap\tilde G_{\xi}$, salah satu dari $\mu H$,
$\mu (\tilde G_{\xi}\setminus H)$ harus bernilai $0$, menurut apakah
$H^{-1}[\{\xi\}]$ terhitung atau tidak.   Jadi $E\cap G_{\xi}$ dan
$E\cap\tilde G_{\xi}$, jika tidak terabaikan, harus merupakan atom, dan
$E$ harus memuat suatu atom.   Karena $E$ sembarang, $\mu$ atomik murni.\ \Qed}

\cmmnt{\medskip

{\bf (h)} Jadi $(X,\Sigma,\mu)$ lengkap, ditentukan secara lokal, dan
atomik murni, tetapi tidak dapat dilokalkan.
}%end of comment

\ifdim\pagewidth>467pt\fontdimen3\tenbf=3pt\fontdimen4\tenbf=2pt\fi
\leader{*216E}{}{\bf Ruang lengkap, ditentukan secara lokal, dan dapat
dilokalkan yang tidak dapat dilokalkan secara ketat}\cmmnt{ Untuk
konstruksi terakhir dan yang paling menarik, kita memerlukan suatu hasil
tak sepele dalam kombinatorika infinit, yang telah saya tuliskan dalam
2A1P:  jika $I$ sembarang himpunan, dan
$\langle f_{\alpha}\rangle_{\alpha\in A}$ suatu keluarga dalam
$\{0,1\}^I$, himpunan fungsi dari $I$ ke $\{0,1\}$, dengan $\#(A)$ lebih
besar secara ketat daripada $\frak c$, kardinalitas kontinuum, dan jika
$\langle K_{\alpha}\rangle_{\alpha\in A}$ sembarang keluarga himpunan
bagian terhitung dari $I$, maka harus ada $\alpha$, $\beta\in A$ yang
berbeda sedemikian sehingga $f_{\alpha}$ dan $f_{\beta}$ bersepakat pada
$K_{\alpha}\cap K_{\beta}$.
\fontdimen3\tenbf=1.92pt
\fontdimen4\tenbf=1.28pt

Dengan berbekal fakta ini, saya melanjutkan sebagai berikut.

\medskip

}{\bf (a)} Misalkan $C$ sembarang himpunan dengan kardinalitas lebih
besar daripada $\frak c$.   Tetapkan $I=\Cal PC$ dan $X=\{0,1\}^I$.
Untuk $\gamma\in C$, definisikan $x_{\gamma}\in X$ dengan menetapkan
bahwa $x_{\gamma}(\Gamma)=1$ jika $\gamma\in\Gamma\subseteq C$ dan
$x_{\gamma}(\Gamma)=0$ jika $\gamma\notin\Gamma\subseteq C$.   Misalkan
$\Cal
K$ keluarga himpunan bagian terhitung dari $I$, dan untuk
$K\in\Cal K$, $\gamma\in C$ tetapkan

\Centerline{$F_{\gamma K}=\{x:x\in X,\,x\restr K=x_{\gamma}\restr
K\}\subseteq X$.}

\noindent Misalkan

$$\eqalign{\Sigma_{\gamma}=\{E:E\subseteq X,
&\text{ terdapat }K\in\Cal K\text{ sedemikian sehingga }F_{\gamma
K}\subseteq E\cr
&\text{ atau terdapat }K\in\Cal K\text{ sedemikian sehingga }F_{\gamma K}\subseteq
X\setminus E\}.\cr}$$

\noindent Maka $\Sigma_{\gamma}$ merupakan aljabar-sigma dari
himpunan-himpunan bagian $X$.   \prooflet{\Prf\ (i)
$F_{\gamma\emptyset}\subseteq
X\setminus\emptyset$ sehingga
$\emptyset\in\Sigma_{\gamma}$.   (ii) Definisi $\Sigma_{\gamma}$ simetris
antara $E$ dan $X\setminus E$, sehingga
$X\setminus E\in\Sigma_{\gamma}$ setiap kali $E\in\Sigma_{\gamma}$.
(iii) Misalkan $\sequencen{E_n}$ suatu barisan dalam $\Sigma_{\gamma}$,
dengan gabungan $E$.   ($\alpha$) Jika terdapat $n\in\Bbb N$,
$K\in\Cal K$ sedemikian sehingga $F_{\gamma K}\subseteq E_n$, maka
$F_{\gamma K}\subseteq E$, sehingga $E\in\Sigma_{\gamma}$.
($\beta$) Jika tidak, untuk setiap $n\in\Bbb N$ terdapat $K_n\in\Cal K$
sedemikian sehingga $F_{\gamma,K_n}\subseteq
X\setminus E_n$.   Tetapkan
$K=\bigcup_{n\in\Bbb N}K_n\in\Cal K$.   Maka

$$\eqalign{F_{\gamma K}
&=\{x:x\restr K=x_{\gamma}\restr K\}
=\{x:x\restr K_n=x_{\gamma}\restr K_n\text{ untuk setiap }n\in\Bbb N\}\cr
&=\bigcap_{n\in\Bbb N}F_{\gamma,K_n}
\subseteq\bigcap_{n\in\Bbb N}X\setminus E_n
=X\setminus E,\cr}$$

\noindent sehingga kembali $E\in\Sigma_{\gamma}$.
Karena $\sequencen{E_n}$ sembarang, $\Sigma_{\gamma}$ merupakan
aljabar-sigma.\ \Qed}

\medskip

{\bf (b)} Tetapkan

\Centerline{$\Sigma=\bigcap_{\gamma\in C}\Sigma_{\gamma}$;}

\noindent maka $\Sigma$\cmmnt{, sebagai irisan aljabar-aljabar-sigma,}
merupakan aljabar-sigma dari himpunan-himpunan bagian $X$
\cmmnt{ (lihat 111Ga)}. Definisikan $\mu:\Sigma\to[0,\infty]$ dengan
menetapkan

$$\eqalign{\mu E&=\#(\{\gamma:x_{\gamma}\in E\})
\text{ jika himpunan ini hingga},\cr
&=\infty\text{ jika tidak};\cr}$$

\noindent maka $\mu$ merupakan suatu ukuran.

\medskip

{\bf (c)} Kelak akan berguna untuk mengetahui sesuatu tentang
himpunan-himpunan

\Centerline{$G_D=\{x:x\in X,\,x(D)=1\}$}

\noindent untuk $D\subseteq C$.   Khususnya, setiap $G_D$ termasuk dalam
$\Sigma$.   \prooflet{\Prf\  Jika $\gamma\in D$, maka
$x_{\gamma}(D)=1$ sehingga $G_D=F_{\gamma,\{D\}}\in\Sigma_{\gamma}$.
Jika $\gamma\in C\setminus D$, maka $x_{\gamma}(D)=0$ sehingga
$G_D=X\setminus
F_{\gamma,\{D\}}\in\Sigma_{\gamma}$.\ \QeD}
Selain itu\cmmnt{, tentu saja,} $\{\gamma:x_{\gamma}\in G_D\}=D$.

\medskip

{\bf (d)} $(X,\Sigma,\mu)$ lengkap.   \prooflet{\Prf\ Andaikan
$A\subseteq
E\in\Sigma$ dan $\mu E=0$.   Untuk setiap
$\gamma\in C$, $E\in\Sigma_{\gamma}$ dan $x_{\gamma}\notin E$, sehingga
$F_{\gamma
K}\not\subseteq E$ untuk setiap $K\in\Cal K$ dan terdapat
$K\in\Cal K$ sedemikian sehingga
\Centerline{$F_{\gamma K}\subseteq X\setminus E\subseteq X\setminus A$.}

\noindent Jadi $A\in\Sigma_{\gamma}$;  karena $\gamma$ sembarang,
$A\in\Sigma$.   Karena $A$ sembarang, $\mu$ lengkap.\ \Qed}

\medskip

{\bf (e)} $(X,\Sigma,\mu)$ semihingga.   \prooflet{\Prf\ Misalkan
$E\in\Sigma$ suatu himpunan berukuran positif.   Maka harus ada suatu
$\gamma\in C$ sedemikian sehingga $x_{\gamma}\in E$.   Tinjau
$E'=E\cap G_{\{\gamma\}}$.   Karena $x_{\gamma}\in E'$, maka
$\mu E'\ge 1>0$.   Di lain pihak,
$\mu
G_{\{\gamma\}}=\#(\{\delta:\delta\in\{\gamma\}\})=1$, sehingga
$\mu E'=1$. Karena $E$ sembarang, $\mu$ semihingga.\ \Qed}

\medskip

{\bf (f)} $(X,\Sigma,\mu)$ dapat dilokalkan.   \prooflet{\Prf\ Misalkan
$\Cal E$ sembarang subhimpunan dari $\Sigma$.   Tetapkan
$D=\{\delta:\delta\in
C,\,x_{\delta}\in\bigcup\Cal E\}$.   Tinjau
$G_D$. Untuk $H\in\Sigma$,

$$\eqalign{\mu(E\setminus H)&=0\text{ untuk setiap }E\in\Cal E\cr
&\iff x_{\gamma}\notin E\setminus H\text{ untuk setiap }E\in\Cal E,
\gamma\in C\cr
&\iff x_{\gamma}\in H\text{ untuk setiap }\gamma\in D\cr
&\iff x_{\gamma}\notin G_D\setminus H\text{ untuk setiap }\gamma\in C\cr
&\iff \mu(G_D\setminus H)=0.\cr}$$

\noindent Jadi $G_D$ merupakan supremum esensial bagi $\Cal E$ dalam
$\Sigma$.   Karena $\Cal E$ sembarang, $\mu$ dapat dilokalkan.\ \Qed}

\medskip

{\bf (g)} $(X,\Sigma,\mu)$ tidak dapat dilokalkan secara ketat.
\prooflet{\Prf\Quer\
Andaikan, jika mungkin, bahwa $\langle X_j\rangle_{j\in J}$ merupakan
suatu dekomposisi dari $(X,\Sigma,\mu)$.   Tetapkan
$J'=\{j:j\in J,\,\mu X_j>0\}$. Untuk setiap $j\in J'$, himpunan
$C_j=\{\gamma:x_{\gamma}\in X_j\}$ harus hingga dan tak kosong.   Selain
itu, untuk setiap $\gamma\in C$, harus ada suatu $j\in J$ sedemikian
sehingga $\mu(G_{\{\gamma\}}\cap X_j)>0$, dan dalam kasus ini $j\in J'$
serta $\gamma\in C_j$.   Jadi $C=\bigcup_{j\in J'}C_j$. Karena
$\#(C)>\frak c$, $\#(J')>\frak c$ (2A1Ld).

Untuk setiap $j\in J'$, pilih $\gamma_j\in C_j$.   Maka

\Centerline{$x_{\gamma_j}\in X_j\in\Sigma\subseteq\Sigma_{\gamma_j}$,}

\noindent sehingga harus ada suatu $K_j\in\Cal K$ sedemikian sehingga
$F_{\gamma_j,K_j}\subseteq X_j$.

Pada titik ini saya akhirnya menggunakan hasil yang dikutip pada awal
contoh ini.   Karena $\#(J')>\frak c$, harus ada $j$, $k\in
J'$ yang
berbeda sedemikian sehingga $x_{\gamma_j}$ dan $x_{\gamma_k}$ bersepakat
pada $K_j\cap K_k$. Karena itu kita dapat mendefinisikan $x\in X$ dengan
menetapkan bahwa

$$\eqalign{x(\delta)&=x_{\gamma_j}(\delta)\text{ jika }\delta\in K_j,\cr
&=x_{\gamma_k}(\delta)\text{ jika }\delta \in K_k,\cr
&=0\text{ jika }\delta\in I\setminus(K_j\cup K_k).\cr}$$

\noindent Sekarang

\Centerline{$x\in F_{\gamma_j,K_j}\cap F_{\gamma_k,K_k}
\subseteq X_j\cap X_k$,}

\noindent dan $X_j\cap X_k\ne\emptyset$;  ini bertentangan dengan
asumsi bahwa $X_j$ membentuk suatu dekomposisi dari $X$.\ \Bang\Qed}

\medskip

{\bf (h)} $(X,\Sigma,\mu)$ atomik murni.   \prooflet{\Prf\ Jika
$E\in\Sigma$ dan $\mu E>0$, maka (sebagaimana dicatat dalam (e) di atas)
terdapat $\gamma\in C$ sedemikian sehingga
$\mu(E\cap G_{\{\gamma\}})=1$;  kini $E\cap G_{\{\gamma\}}$ harus
merupakan suatu atom.\ \Qed}

\cmmnt{\medskip

{\bf (i)} Dengan demikian $(X,\Sigma,\mu)$ merupakan ruang ukur lengkap,
ditentukan secara lokal, dapat dilokalkan, dan atomik murni yang tidak
dapat dilokalkan secara ketat.
}%end of comment

\exercises{
\leader{216X}{Latihan dasar (a)}
%\spheader 216Xa
Dalam konstruksi 216C, tunjukkan bahwa ukuran subruang pada
$\{1\}\times I$ tidak semihingga.
%216C

\spheader 216Xb Andaikan, dalam 216D, bahwa $I=\omega_1$.   (i) Tunjukkan
bahwa himpunan $\{(\xi,\eta):\xi\le\eta<\omega_1\}$ terukur oleh ukuran
yang dikonstruksi dengan metode \Caratheodory dari
$\mu^*\restrp\Cal P(I\times I)$, tetapi tidak oleh ukuran subruang pada
$I\times I$.   (ii) Dari sini, atau dengan cara lain, tunjukkan bahwa
ukuran subruang pada $I\times I$ tidak ditentukan secara lokal.
%216D

\spheader 216Xc Dalam 216Ya, 252Yq, dan 252Ys di bawah, saya menunjukkan
cara mengonstruksi versi tanpa atom dari 216C, 216D, dan 216E, yakni ruang
ukur lengkap tanpa atom yang pertama dapat dilokalkan tetapi tidak
ditentukan secara lokal, yang kedua ditentukan secara lokal tetapi tidak
dapat dilokalkan, dan yang ketiga ditentukan secara lokal dan dapat
dilokalkan tetapi tidak dapat dilokalkan secara ketat.   Tunjukkan
bagaimana jumlah langsung dari ruang-ruang ini, bersama dengan ukuran
pencacahan dan contoh-contoh yang dijelaskan dalam bab ini, dapat disusun
untuk memberikan semua 56 contoh yang dituntut oleh pembahasan dalam
pengantar bagian ini.
%216E

\leader{216Y}{Latihan lanjutan (a)}
%\spheader 216Ya
Misalkan $\lambda$ ukuran Lebesgue pada $[0,1]$, dan $\Lambda$ domainnya.
Tetapkan $Y=[0,1]\times\{0,1\}$ dan tuliskan

\Centerline{$\Tau=\{F:F\subseteq Y,\,F^{-1}[\{0\}]\in\Lambda\}$,}

\Centerline{$\nu F=\lambda F^{-1}[\{0\}]$ untuk setiap $F\in\Tau$.}

\noindent Tetapkan

\Centerline{$\Tau_0=\{F:F\in\Tau,\,F^{-1}[\{0\}]\symmdiff
F^{-1}[\{1\}]$ bersifat $\lambda$-terabaikan$\}$.}

\noindent Misalkan $I$ suatu himpunan tak terhitung.   Tetapkan
$X=Y\times I$,

\Centerline{$\Sigma=\{E:E\subseteq X,\,E^{-1}[\{i\}]\in\Tau$ untuk setiap $i\in
I$, $\{i:E^{-1}[\{i\}]\notin \Tau_0\}$ terhitung$\}$,}

\Centerline{$\mu E=\sum_{i\in I}\nu E^{-1}[\{i\}]$ untuk $E\in\Sigma$.}

\noindent (i) Tunjukkan bahwa $(Y,\Tau,\nu)$ dan
$(Y,\Tau_0,\nu\restrp\Tau_0)$ merupakan ruang peluang lengkap, dan bahwa
untuk setiap $F\in\Tau$ terdapat $F'\in\Tau_0$ sedemikian sehingga
$\nu(F\symmdiff F')=0$.   (ii) Tunjukkan bahwa $(X,\Sigma,\mu)$ merupakan
ruang ukur lengkap tanpa atom yang dapat dilokalkan tetapi tidak
ditentukan secara lokal.
%216C

\spheader 216Yb Definisikan suatu ukuran $\mu$ pada
$X=\omega_2\times\omega_2$ sebagai berikut.   Ambil $\Sigma$ sebagai
aljabar-sigma dari himpunan-himpunan bagian $X$ yang dibangkitkan oleh

\Centerline{$\{A\times\omega_2:A\subseteq\omega_2\}
\cup\{\omega_2\times\alpha:\alpha<\omega_2\}$.}

\noindent Untuk $E\in\Sigma$ tetapkan

\Centerline{$W(E)=\{\xi:\xi<\omega_2$, $\sup E[\{\xi\}]=\omega_2\}$,}

\noindent dan tetapkan $\mu E=\#(W(E))$ jika ini hingga, $\infty$ jika
tidak. Tunjukkan bahwa $\mu$ merupakan ukuran pada $X$, dapat dilokalkan,
dan ditentukan secara lokal, tetapi tidak memiliki himpunan terabaikan
yang ditentukan secara lokal.   Carilah suatu subruang $Y$ dari $X$
sedemikian sehingga ukuran subruang pada $Y$ tidak semihingga.
%216D

\spheader 216Yc Tunjukkan bahwa dalam ruang yang dijelaskan dalam 216E,
setiap himpunan memiliki selubung terukur, tetapi hal ini tidak benar
dalam ruang-ruang 216C dan 216D.
%216E

\spheader 216Yd Tetapkan $X=\omega_1\times\omega_2$.   Untuk
$E\subseteq X$ tetapkan

\Centerline{$A(E)=\{\zeta:$ untuk suatu $\xi$, tepat satu dari
$(\xi,\zeta)$, $(\xi,\zeta+1)$ termasuk dalam $E\}$,}

\ifdim\pagewidth>400pt
\Centerline{$B(E)=\{\zeta:$ terdapat $\xi$, $\zeta'$ sedemikian sehingga
$\zeta<\zeta'<\omega_2$ dan tepat satu dari
$(\xi,\zeta)$, $(\xi,\zeta')$ termasuk dalam $E\}$,}
\else
$$\eqalign{B(E)
&=\{\zeta:\text{ terdapat }\xi,\,\zeta'\text{ sedemikian sehingga }
\zeta<\zeta'<\omega_2\cr&
\mskip100mu\text{ dan tepat satu dari }
(\xi,\zeta),\,(\xi,\zeta')\text{ termasuk dalam }E\},
\cr}$$
\fi

\Centerline{$W(E)=\{\xi:\#(E[\{\xi\}])=\omega_2\}$.}

\noindent Misalkan $\Sigma$ adalah himpunan dari himpunan-himpunan bagian
$E$ dari $X$ sedemikian sehingga $A(E)$ terhitung dan
$\#(B(E))\le\omega_1$.   Untuk $E\in\Sigma$, tetapkan
$\mu E=\#(W(E))$ jika ini hingga, $\infty$ jika tidak.   (i) Tunjukkan
bahwa $(X,\Sigma,\mu)$ merupakan ruang ukur.   (ii) Tunjukkan bahwa jika
$\hat\mu$ adalah pelengkapan dari $\mu$, maka domainnya adalah himpunan
dari himpunan-himpunan bagian $E$ dari $X$ sedemikian sehingga $A(E)$
terhitung, dan $\hat\mu$ dapat dilokalkan secara ketat.   (iii) Tunjukkan
bahwa $\mu$ tidak dapat dilokalkan secara ketat.
%216Yb 216D

\spheader 216Ye\dvAnew{2015} Tunjukkan bahwa terdapat suatu ruang ukur
lengkap, tanpa atom, dan semihingga dengan suatu subhimpunan satu anggota
yang tidak terabaikan.
\Hint{tetapkan $X=(\omega_1\times[0,1])\cup\{\omega_1\}$ dan misalkan
$\Sigma$ aljabar-sigma dari himpunan-himpunan bagian $X$ yang
dibangkitkan oleh $\{\{\xi\}\times E:\xi<\omega_1$,
$E\subseteq[0,1]$ terukur Lebesgue$\}$}.
%215E

}%end of exercises

\endnotes{
\Notesheader{216} Contoh-contoh 216C-216E %216C 216D 216E
dirancang untuk, bersama dengan ukuran Lebesgue, membentuk suatu landasan
bagi konstruksi seperangkat contoh lengkap untuk konsep-konsep yang
didaftar dalam 211A-211K.   %211A 211B 211C 211D 211E 211F 211G 211H 211I 211J 211K
Kita tidak benar-benar mengharapkan untuk menjumpai fenomena-fenomena ini
dalam penerapan, tetapi pemahaman yang jelas tentang kemungkinan yang
diperlihatkan oleh contoh-contoh ini merupakan bagian dari penilaian yang
tepat terhadap kelangkaannya.   Tentu saja, jika kita menambahkan sifat
lain pada daftar kita -- misalnya, sifat mempunyai himpunan terabaikan
yang ditentukan secara lokal (213I), atau sifat bahwa setiap himpunan
bagian harus mempunyai selubung terukur (213Xl) -- maka terdapat hasil
positif lebih lanjut untuk melengkapi 211L, dan lebih banyak contoh untuk
dicari, seperti 216Yb. Tetapi sudah waktunya, bahkan mungkin sudah lewat
waktunya, bagi kita untuk kembali kepada teorema-teorema klasik yang
berlaku pada ruang-ruang ukur di pusat bidang ini.
}%end of comment

\discrpage




